\section{ADR-004: Validated types at every boundary}
\label{adr:004}

\subsection*{Status}

Accepted; in force at revision \texttt{\projectrevisionshort}.

\subsection*{Context}

Data enters from CSV files, HTTP request bodies, JSONL streams, JSON files on
disk, and language-model responses. All of these are untrusted and
loosely-shaped.

\subsection*{Decision}

Every value crossing a boundary is a Pydantic model, and constraints are
declared on the fields rather than checked imperatively: rating bounds,
minimum text length, embedding-dimension range, metric ranges in $[0,1]$,
non-negative counters, and cross-field rules such as
\texttt{min\_rating} $\leq$ \texttt{max\_rating}. Persisted state --- stored
conversations, jobs, reports, human feedback, active-learning state --- is
re-validated on load rather than trusted.

\subsection*{Consequences}

\begin{itemize}
  \item Invalid input fails at the boundary with a message naming the field, so
    interior code can assume well-formed data and needs no defensive checks.
  \item FastAPI derives request validation and the OpenAPI schema from the same
    models, so the HTTP contract and the domain contract cannot diverge.
  \item \texttt{schemas} acquires the highest fan-in in the package (26 of 43
    modules import it), which the coupling table shows. The shared vocabulary
    is a deliberate coupling.
  \item Pydantic reaches the foundation layer and therefore, transitively,
    every layer. Replacing it would be a repository-wide change.
  \item Model construction costs more than a dictionary. The repository accepts
    this; the benchmark harness measures the pipeline including that cost.
\end{itemize}

\subsection*{Alternatives}

Dataclasses with hand-written validation, or dictionaries with runtime
assertions. Both would require duplicating the HTTP schema separately from the
domain types.

\subsection*{Evidence}

\srcfile{src/feedback\_intelligence\_agent/schemas.py},
\srcfile{data\_contracts.py}, \srcfile{config.py}, \srcfile{evaluation.py},
\srcfile{memory.py}, \srcfile{jobs.py}, \srcfile{reports.py};
\srcfile{AGENTS.md} (``Use Pydantic models for external request and response
schemas''; ``Raise specific exceptions for data quality issues'');
\srcfile{pyproject.toml} (\texttt{plugins = ["pydantic.mypy"]}).
